This approach is also referred to as the `flat histogram sampling'. In Wang-Landau sampling [\cref{fig:sampling_methods}(f)],\cite{wang2001efficient,wang2001determining} a density-of-states histogram is being built over a targeted range of energies $[E_{\min},E_{\max} ]$ as the sampling progresses. When a new sample with energy of $E_{new}$ is generated, the histogram at this energy state will be updated and the `flatness' of the histogram will be assessed. The new sample will be accepted if it leads to a more flattened histogram and rejected otherwise. Provided that $E_{\max}  - E_{\min}$ is (much) larger than the thermal energy $k_B T$, for which a canonical MC sampling will be performed otherwise, the Wang-Landau acceptance criteria will, therefore, encourage high-energy configurations to be accepted as the intermediate steps (anchoring points) for crossing the high energy barriers in the hope of reaching other deep local minima. Essentially, the flat histogram can be regarded as the result of combining the Boltzmann distributions from multiple temperatures, therefore, the Wang-Landau sampling leads to a \textit{multicannonical ensemble} of configurations for a given material.